% generated by figures-src/make_tables.py -- do not edit
\begin{table}[tb]
\centering
\caption{S\=amavedic notation: every one of the 1{,}844 \=Arcika verses carries
exactly one typed disposition, so a reader can never mistake a withheld verse for a
verse the tradition does not notate. The third withheld class did not exist until an
adversarial gate falsified the stated reasons of 39 rows.}
\label{tab:samaveda}
\small
\begin{tabularx}{\linewidth}{@{}lr>{\raggedright\arraybackslash}X@{}}
\toprule
Disposition & Verses & \\
\midrule
notation present, source-explicit & 1{,}136 & \footnotesize the witness prints marks on this verse \\
\quad withheld: near line, witness disagreement & 565 & \footnotesize needs philological adjudication \\
\quad withheld: no near line, probable absence & 104 & \footnotesize the witness may not carry the verse \\
\quad withheld: occurrence consumed by a repeat & 39 & \footnotesize \textbf{added by Gate C} (\cref{sec:case-samaveda}) \\
\midrule
\textbf{Total} & \textbf{1{,}844}
& \footnotesize verses with no disposition at all: \textbf{0} \\
\bottomrule
\end{tabularx}
\end{table}
